\documentclass[14pt,a4paper]{extarticle}
\usepackage[T2A]{fontenc}\usepackage[utf8]{inputenc}\usepackage[russian]{babel}
\usepackage{geometry}\geometry{left=30mm,right=15mm,top=20mm,bottom=20mm}
\setlength{\parindent}{1.25cm}\setlength{\parskip}{0.3em}\sloppy
\begin{document}\begin{titlepage}\begin{center}
САНКТ-ПЕТЕРБУРГСКИЙ НАЦИОНАЛЬНЫЙ\\ ИССЛЕДОВАТЕЛЬСКИЙ УНИВЕРСИТЕТ ИТМО
\vspace{1.5cm}

Дисциплина: Бэк-энд разработка
\vspace{1.5cm}

\textbf{Отчёт}\\[0.5cm]
Домашняя работа 8\\[0.5cm]
\textbf{Настройка CI/CD для облачной инфраструктуры}
\vfill
\begin{flushright}Выполнил:\\ Цатинян Артём\\ Группа БР1.2\\[0.8cm] Проверил:\\ Добряков Д. И.\end{flushright}
\vfill Санкт-Петербург\\ 2026 г.\end{center}\end{titlepage}
\section{Задача и состояние}
Описать автоматизированную интеграцию и доставку микросервисов в облачную инфраструктуру. Ниже представлен проект пайплайна; его выполнение в облаке не подтверждено.
\section{Исходные условия}
В приложении используются сервисы identity, catalog, booking и review на Kotlin и Spring Boot. Целевой стенд предложен в ДЗ7; облачное развёртывание из ЛР7 ещё требуется выполнить.
\section{Непрерывная интеграция}
При создании pull request предлагается запускать Gradle-сборку и доступные тесты, проверять конфигурацию, собирать образы изменённых сервисов и сканировать их на уязвимости. Версии действий и зависимости фиксируются; результат должен быть виден в проверках pull request.
\section{Доставка}
После слияния в целевую ветку предлагается публиковать образы с тегом коммита, обновлять версии образов в Helm-чарте и выполнять rollout в тестовое окружение. После проверки доступности можно согласовать продвижение в рабочее окружение. Для доступа к реестру и кластеру используются ограниченные права и секреты системы CI.
\section{Контроль и откат}
Проверяются успешность сборки, готовность подов, ответы HTTP и работоспособность зависимостей. При неудачном rollout используется последняя проверенная версия; миграции БД планируются с учётом совместимости. Реальный запуск пайплайна и время восстановления необходимо подтвердить экспериментально.
\section{Вывод}
Описана схема CI/CD; конфигурацию, запуск и подтверждение результатов ещё необходимо выполнить.
\end{document}
